Bei konstanter Geschwindigkeit muss der Motor die Luftwiderstandskraft überwinden (FoTaBe S.~165):
\begin{align}
 F_\mathrm{L} &= \frac{1}{2}c_w\,A\,\rho\,v^2
\end{align}
Die Leistung ist Kraft mal Geschwindigkeit:
\begin{align}
 P &= F_\mathrm{L}\,v = \frac{1}{2}c_w\,A\,\rho\,v^3
\end{align}
Die Leistung ist proportional zur dritten Potenz der Geschwindigkeit. Mit $v_1 = \vlaO$ und $v_2 = \vlbO$ ist
\begin{align}
 \frac{P_2}{P_1} &= \fakF = \left(\frac{\vlbO}{\vlaO}\right)^3 = \fak \approx \result{\fakP{0}{3}}
\end{align}
Eine um einen Zehntel höhere Geschwindigkeit erfordert also eine um rund einen Drittel grössere Leistung.
